\documentclass[10pt,a4paper]{amsart}
\usepackage[margin=19mm]{geometry}
\usepackage{amsmath,amssymb,mathtools}
\usepackage[scr=rsfs,cal=euler]{mathalfa}
\usepackage{xfrac}
\usepackage[unicode,hidelinks,bookmarks=false]{hyperref}
\newcommand{\ie}{i.e.\ }
\newcommand{\cf}{cf.\ }
\newcommand{\resp}{respectively\ }
\newcommand{\Subsection}{\subsection}
\providecommand{\nobreakdash}{\nobreak\hbox{-}}
\setlength{\emergencystretch}{2em}
\multlinegap0pt
\usepackage{bm}
\usepackage{bbm}
\usepackage{dsfont}
\usepackage{xspace}
\usepackage{cancel}
\usepackage{mathtools}
\usepackage{amsthm}
\makeatletter
\@ifundefined{theorem}{\newtheorem{theorem}{Theorem}}{}
\@ifundefined{proposition}{\newtheorem{proposition}[theorem]{Proposition}}{}
\makeatother
\title[A note on the differential smoothness of skew PBW extensions]{A note on the differential smoothness of skew PBW extensions}
\author{Andr\'es Rubiano, Armando Reyes}
\date{Source: arXiv:2410.13967v2 (CC BY 4.0)}
\begin{document}
\maketitle

\noindent\emph{Source and licence.} A note on the differential smoothness of skew PBW extensions. Version arXiv:2410.13967v2 (CC BY 4.0), distributed under the Creative Commons Attribution 4.0 International licence, as recorded in the arXiv licence metadata for this exact version and shown on the abstract page https://arxiv.org/abs/2410.13967v2. Full rights notice: licenses/arxiv-2410.13967-CC-BY-4.0.txt.\par\smallskip

\noindent\textit{Editorial scope.} Counted unit: the Proposition whose source label is \texttt{extaautoSPBW1} in the article (source0.tex lines 407--409, source label \texttt{extaautoSPBW1}), delivered together with the article's own complete original proof of it (source0.tex lines 410--494, one \texttt{proof} environment of 85 lines). No mathematics is written, repaired or re-derived: statement and proof are the article's own text line for line. Required context retained uncounted: none. Every definition, display and label of those lines is retained unchanged. Omitted from this package and not redistributed: 1--406, 495--994. Nothing inside a retained excerpt is omitted or altered: every retained body is delivered exactly as the article's lines are written, so no omission receipt of any kind accompanies this package. Citations: the retained excerpts carry no citation command at all, so no bibliography entry of the article is transcribed and no reference is redistributed. Reference adaptation: none was needed and none was made; every reference inside the delivered unit resolves inside it, so no unresolved reference or citation is produced. Printed slips kept literal: no printed word, symbol, hypothesis or number of the counted statement or of its proof was altered. Layout and class: the article's own class is \texttt{amsart} with packages including amsthm, amssymb, amsfonts, latexsym, mathtools, thmtools, amsmath, lscape, fontenc, tikz-cd, enumitem, longtable, mathrsfs, hyperref; the wrapper is the lane's pinned \texttt{amsart} editorial wrapper at 10pt on A4, which restates the theorem-like environments the retained text needs and adds the article's own macro definitions that the retained text needs (none), with extra wrapper packages (bm, bbm, dsfont, xspace, cancel, mathtools). The delivered numbering is the unit's own. Assets: no retained span contains a figure environment, a graphics-inclusion command, a PDF-page inclusion, a table or a TikZ picture; no image file is copied into this package, the package declares no asset dependency and no image file is redistributed with it. Licence: version arXiv:2410.13967v2 of this article is distributed under the Creative Commons Attribution 4.0 International licence, as recorded in the arXiv licence metadata for this exact version and shown on the abstract page https://arxiv.org/abs/2410.13967v2. A full rights notice is delivered at licenses/arxiv-2410.13967-CC-BY-4.0.txt. Characters: every retained excerpt is pure ASCII, so the delivered engine, XeLaTeX with its default Latin Modern text font, reports no missing character.
\begin{proposition}\label{extaautoSPBW1}
Let $A = \sigma(R)\langle x \rangle$ be a SPBW extension over $R$ of automorphism type. Consider the automorphism $\widetilde{\sigma}:A \rightarrow A$ given by $\widetilde{\sigma}(x) = x$ and $\widetilde{\sigma}(r)=\sigma(r)$ for each $r\in R$. If  $\sigma \circ \delta = \delta \circ \sigma$, then the map $\widetilde{\delta}:A\rightarrow A$ with $\widetilde{\delta}(f(x))=\widetilde{\delta}(a_0+a_1x+ \cdots +a_nx^n)=\delta(a_0)+\delta(a_1)x+ \cdots + \delta(a_n)x^n$ for all $a_i \in R$, $0\leq i\leq n$ is a $\widetilde{\sigma}$-derivation of $A$.
\end{proposition}

\begin{proof}
    Let $p(x) = \sum\limits_{i=0}^n r_i x^i \in A $. Then
    \begin{align*}
        \widetilde{\sigma}(p(x)) = \widetilde{\sigma} \left(\sum_{i=0}^n r_i x^i\right) = \sum_{i=0}^n \widetilde{\sigma}(r_i x^i) = \sum_{i=0}^n\widetilde{\sigma} (r_i) x^i = \sum_{i=0}^n\sigma(r_i) x^i.
    \end{align*}
    
It is clear that the bijectivity of $\widetilde{\sigma}$ is inherited by $\sigma$.

Now, for $p(x) = r x^k$, we have that
\begin{equation*}
    \widetilde{\delta} (p(x)) = \widetilde{\delta}(r x^k) = \widetilde{\sigma}(r)\widetilde{\delta} (x^k) +x  \widetilde{\delta}(r) x^k = \sigma(r)\widetilde{\delta}(x^k) + \delta(r) x^k.
\end{equation*}

Let $r(x) = r x^i$ and $s(x) = s x^j$ for some $i,j \in \mathbb{N}$. The idea is to satisfy the relation given by
\begin{align*}
    \widetilde{\delta} (r(x) s(x)) = \widetilde{\sigma} (r(x)) \widetilde{\delta}(s(x)) + \widetilde{\delta} (r(x)) s(x).
\end{align*}

With this aim, note that 
\begin{align}
    \widetilde{\delta}(r(x) & s(x)) = \widetilde{\delta}(r x^is x^j) \notag \\ 
    = &\ \widetilde{\delta}\left(r\sum_{k=0}^{i}\sum_{f_k\in C_{i-k,k}}f_k(s) x^{i-k} x^j\right) \notag \\
     = &\ \widetilde{\delta}\left(\sum_{k=0}^{i}\sum_{f_k\in C_{i-k,k}} r f_k(s) x^{i-k+j}\right) \notag \\ 
     = &\ \sum_{k=0}^{i}\sum_{f_k\in C_{i-k,k}}\widetilde{\delta}\left(rf_k(s) x^{i-k+j}\right) \notag \\
     = &\ \sum_{k=0}^{i}\sum_{f_k\in C_{i-k,k}}\left(\sigma(rf_k(s))\widetilde{\delta} (x^{i-k+j}) + \delta(rf_k(s)) x^{i-k+j}\right) \notag \\
     = &\ \sum_{k=0}^{i}\sum_{f_k\in C_{i-k,k}}\left(\sigma(r)\sigma(f_k(s))\widetilde{\delta}(x^{i-k+j})+\left(\sigma(r)\delta(f_k(s))+\delta(r)f_k(s)\right) x^{i-k+j}\right) \notag 
     \end{align}
    \begin{align} 
     = &\ \sum_{k=0}^{i} \sum_{f_k\in C_{i-k,k}} \left(\sigma(r)\sigma(f_k(s))\left(x^{i-k}\widetilde{\delta}(x^j) + \widetilde{\delta}(x^{i-k}) x^j\right) \right. \notag \\
     &\ \ + \left. \left(\sigma(r)\delta(f_k(s)) + \delta(r)f_k(s)\right) x^{i-k+j} \right). \label{First24}
\end{align}

On the other hand,
{\small{
\begin{align}
    \widetilde{\sigma}(r(x)) \widetilde{\delta}(s(x)) + \widetilde{\delta}(r(x))s(x) = &\ \widetilde{\sigma}(r x^i)\widetilde{\delta}(s x^j)+\widetilde{\delta}(r x^i) sx^j \notag \\
    = &\ \sigma(r) x^i\left(\sigma(s)\widetilde{\delta}(x^j) + \delta(s) x^j\right) + \left(\sigma(r)\widetilde{\delta}(x^i) + \delta(r) x^i\right)s x^j \notag \\
    = &\ \sigma(r) x^i \sigma(s)\widetilde{\delta}(x^j) + \sigma(r) x^i\delta(s) x^j \\
    &\ + \sigma(r)\widetilde{\delta}(x^i) sx^j + \delta(r) x^is x^j \notag \\
    = &\ \sigma(r)\left(\sum_{k=0}^{i}\sum_{f_k\in C_{i-k,k}}f_k(\sigma(s)) x^{i-k}\right)\widetilde{\delta}(x^j) \\
    &\ + \sigma(r)\left(\sum_{k=0}^{i}\sum_{f_k\in C_{i-k,k}}f_k(\delta(s)) x^{i-k}\right) x^j \notag \\ &\ \ + \sigma(r)\widetilde{\delta}(x^i) s x^j + \delta(r) \left(\sum_{k=0}^{i}\sum_{f_k\in C_{i-k,k}}f_k(s) x^{i-k}\right)x^j \notag \\
    = &\ \sum_{k=0}^{i}\sum_{f_k\in C_{i-k,k}}\left(\sigma(r)f_k(\sigma(s)) x^{i-k}\widetilde{\delta}(x^j) + \sigma(r)f_k(\delta(s)) x^{i-k+j} \right. \\
    &\ + \left. \delta(r)f_k(s) x^{i-k+j}\right) + \sigma(r)\widetilde{\delta}(x^i)s x^j. \label{Second24}
\end{align}
}}

If we compare (\ref{First24}) and (\ref{Second24}), we get that
\begin{align*}
&\ \sum_{k=0}^{i}\sum_{f_k\in C_{i-k,k}}\left((\sigma(f_k(s))-f_k(\sigma(s))) x^{i-k}\widetilde{\delta}(x^j) + (\delta(f_k(s))-f_k(\delta(s))) x^{i-k+j} \right. \\
&\ \ + \left. \sigma(f_k(s))\widetilde{\delta}(x^{i-k}) x^j\right) - \widetilde{\delta}(x^i)s x^j=0.
\end{align*}

Let $\widetilde{\delta}(x^i) = \sum\limits_{l=0}^{n} a_lx^l$, $\widetilde{\delta}(x^j) = \sum\limits_{l=0}^{n}b_l x^l$ and $\widetilde{\delta}(x^{i-k}) = \sum\limits_{l=0}^{n}c_l x^l$. Then,
\begin{align*}
&\ \notag \sum_{k=0}^{i}\sum_{f_k\in C_{i-k,k}}\left((\sigma(f_k(s))-f_k(\sigma(s))) x^{i-k}\sum_{l=0}^{n}b_l x^l+(\delta(f_k(s))-f_k(\delta(s))) x^{i-k+j} \right. \\ \notag
&\ \ \left. + \ \sigma(f_k(s))\sum_{l=0}^{n}c_l x^l x^j\right) - \sum_{l=0}^{n}a_l x^ls x^j=0. \notag
\end{align*}

\begin{align*}
&\ \sum_{k=0}^{i}\sum_{f_k\in C_{i-k,k}}\left((\sigma(f_k(s))-f_k(\sigma(s)))\sum_{l=0}^{n}\left(\sum_{p=0}^{i-k}\sum_{h_p\in C_{i-k-p,p}}h_p(b_l) x^{i-k-p}\right) x^l \right. \\ \notag
&\ \ \left. + \ (\delta(f_k(s))-f_k(\delta(s))) x^{i-k+j} +\sigma(f_k(s))\sum_{l=0}^{n}c_l x^l x^j\right)\\ \notag
&\ -\sum_{l=0}^{n}a_l\left(\sum_{k=0}^l\sum_{g_k\in C_{l-k,k}}g_k(s) x^{l-k}\right) x^j = 0. 
\end{align*}

\begin{align}\label{Eq1}
&\ \sum_{k=0}^{i}\sum_{f_k\in C_{i-k,k}}\left((\sigma(f_k(s))-f_k(\sigma(s)))\sum_{l=0}^{n}\sum_{p=0}^{i-k}\sum_{h_p\in C_{i-k-p,p}}h_p(b_l)x^{i-k-p+l} \right. \\ \notag
&\ \ \left. + \ (\delta(f_k(s))-f_k(\delta(s)))x^{i-k+j} +\sigma(f_k(s))\sum_{l=0}^{n}c_l x^{l+j}\right)\\ \notag
&\ \ - \ \sum_{l=0}^{n}\sum_{k=0}^l\sum_{g_k\in C_{l-k,k}}a_lg_k(s)x^{l-k+j} = 0. 
\end{align}

When $k=0$, we put attention on the monomial $\sum_{f_0\in C_{i,0}}(\delta(f_0(s))-f_0(\delta(s)))x^{i+j}$. This is the monomial of the highest degree and since it is equal to zero, it follows that $\sum_{f_0\in C_{i,0}} \left(\delta(f_0(s))-f_0(\delta(s))\right) = 0$. Since the power $i$ is arbitrary, it can be seen that $\delta(\sigma(s))=\sigma(\delta(s))$. In the same way, the value of $s$ is arbitrary, which means that $\delta\sigma = \sigma\delta$.

Now, expression {\rm (}\ref{Eq1}{\rm )} can be written as
\begin{align*}
\sum_{k=0}^{i}\sum_{f_k\in C_{i-k,k}}\left(\sigma(f_k(s))\sum_{l=0}^{n}c_lt^{l+j}\right)-\sum_{l=0}^{n}\sum_{k=0}^l\sum_{g_k\in C_{l-k,k}}a_lg_k(s)x^{l-k+j} = &\ 0,
\end{align*}

or equivalently, 
\begin{align*}
\sum_{l=0}^{n}\left(\sum_{k=0}^{i}\sum_{f_k\in C_{i-k,k}}\sigma(f_k(s))c_lx^{l}-\sum_{k=0}^l\sum_{g_k\in C_{l-k,k}}a_lg_k(s)x^{l-k}\right) = &\ 0.
\end{align*}

Since this equality holds for all $s\in R$, any automorphism $\sigma$ and every $\sigma$-derivation $\delta$ of $R$, necessarily $a_l=c_l=0$ for all $0\leq l \leq n$. By changing the values of $i$ with $j$ we obtain that $b_l=0$. Then $\widetilde{\delta}(x)=0$. This means that 
$\widetilde{\delta}(f(x))=\widetilde{\delta}(a_0+a_1x+ \cdots +a_nx^n)=\delta(a_0)+\delta(a_1)x+ \cdots + \delta(a_n)x^n$.
\end{proof}

\end{document}
